\chapter*{Résumé}
\addcontentsline{toc}{chapter}{Résumé}

Ce rapport présente la conception et le développement d’un système automatisé
d’auto-défense réseau, intitulé \textit{Automated Network Self-Defense System}
(ANSDS). Ce projet a été réalisé dans le cadre d’un stage d’été au sein de
l’Agence Nationale de la Cybersécurité (ANCS).

L’objectif principal est de concevoir un système capable d’analyser le trafic
réseau, de détecter automatiquement des comportements suspects, de transmettre
des alertes de sécurité en temps réel et d’appliquer une réponse automatique
sur l’infrastructure réseau.

Une plateforme expérimentale a été développée sous GNS3 en intégrant un routeur
Cisco, un commutateur, un serveur Ubuntu, une machine Kali Linux ainsi qu’un
accès Internet. Cet environnement permet de reproduire différents scénarios
d’activité réseau et de valider progressivement les fonctionnalités du système.

Le moteur de détection a été développé en Python avec la bibliothèque Scapy.
Il permet notamment d’analyser les tentatives de connexion SSH et de détecter
des comportements associés aux scans réseau Nmap. La détection prend en compte
plusieurs indicateurs tels que les signatures TCP, les ports ciblés, la
répétition des paquets et la fenêtre temporelle.

Le système intègre également un mécanisme de corrélation avec l’activité SSH
afin de renforcer la décision de sécurité. Lorsqu’un événement suspect est
confirmé, une notification est automatiquement transmise via Telegram.
Le module Netmiko permet ensuite d’interagir avec le routeur Cisco afin
d’appliquer dynamiquement une ACL de blocage.

Le blocage est temporaire et sa durée est configurée par le système.
À son expiration, l’ACL est automatiquement supprimée afin de rétablir l’état
initial du réseau. Un mécanisme de journalisation permet également de conserver
une trace des principales étapes du traitement.

Les expérimentations réalisées ont permis de valider la détection des scans
Nmap SYN, FIN, NULL et Xmas, la corrélation avec l’activité SSH, l’envoi des
alertes Telegram, l’application du blocage Cisco ainsi que le déblocage
automatique.

\textbf{Mots-clés :} Cybersécurité, ANSDS, SSH, Nmap, Scapy, GNS3, Cisco IOS,
Netmiko, Telegram, détection réseau, auto-défense réseau.